\chapter{The Notebook}\label{ch:notebook}

The notebook evaluates nothing. That is its whole design, and it was decided before a line of it
was written: the frontend owns no printer, no simplifier, no numeric evaluator; it sends source to
an engine and draws what comes back. This chapter is about what that leaves for the notebook to
do, which turns out to be more than expected, and about the three ways the engine reaches it.

\section{One engine, three hosts}

The engine is a Lean program with a JSON-RPC surface (Chapter~\ref{ch:what}). It runs in three
places from one source:

\begin{itemize}[nosep]
  \item \textbf{In the browser}, compiled to WebAssembly and loaded in a web worker. The runtime
    and \lc{Init} are built from source for \texttt{wasm32}, because Lean stopped shipping a
    prebuilt one after v4.15 and the project wants current Lean. The module is 1.8~MB and
    answers its first request 96~ms after the worker starts. No threads, so no cross-origin
    isolation headers: a self-hosted notebook is a static directory.
  \item \textbf{Natively}, behind an HTTP host, for a shared server.
  \item \textbf{On the command line}, over stdio, which is how the tests drive it.
\end{itemize}

The notebook does not know which. It holds an \lc{EngineClient} built from a \lc{Transport},
and a transport moves strings.

\section{What the notebook draws}

Everything the notebook shows is either a string from the engine or a picture described by it.
The engine renders terms to text and to LaTeX, and with paths requested wraps every subterm in an
\texttt{\textbackslash htmlData} annotation; the notebook typesets the LaTeX and attaches a click
handler that reads the path back. That is the whole of selection.

The derivation panel is a list of steps. A step click shows the trail up to it
(Chapter~\ref{ch:origins}); a nested derivation --- the simplification inside a derivative step,
or the differentiation inside an integration check --- opens as a sub-panel with the same
machinery, and the status shown for a step is the weakest status of anything beneath it.

Pictures follow protocol rule 5: a new capability is an optional field, never a change. A plot is
a list of samples with \texttt{null} where the function has no finite value; the notebook draws the
curve, breaks it at the gaps, and trims the tails of the range so an asymptote does not flatten
the rest. A Hasse diagram is a list of elements with heights and a list of covers; the notebook
lays them out in layers. In both cases the engine simplified the formula or decided the order,
recorded a derivation, and described the picture; nothing in the drawing evaluates.

\section{Files, and the session behind them}

A notebook file (\texttt{.chalk}) is JSON: the cells' sources, their saved outputs and derivations, and the studio's
scenes. Opening one shows the saved outputs at once and then re-runs every cell in order, so that
the engine's session --- and with it \texttt{explain}, which needs the derivation to be in the
engine's memory --- matches what is on screen. The notebook also autosaves to the browser after
every run and comes back on reload the same way.

\begin{logbox}[The stale module]
  A browser served last week's WebAssembly for a whole afternoon while the tests passed. Every
  asset URL now carries a build stamp, and the notebook logs the engine's version next to the
  worker's start-up time. Not a finding about mathematics; a finding about trust.
\end{logbox}

\section{Input}

Two small conveniences turned out to matter for the worlds of Chapter~\ref{ch:worlds}. Typing a
backslash opens a completion list of every symbol abbreviation, filtered as you type; at
\texttt{\textbackslash lam} it shows the one entry with the symbol it produces, $\lambda$, on the
right, and Tab or a following space inserts it. The same table covers \texttt{\textbackslash pi},
\texttt{\textbackslash e} (Euler's number, which the engine reads as $\exp 1$ so every rule about
$\exp$ applies to it), \texttt{\textbackslash phi} and the Greek alphabet. This is the input method of the Lean editor
extension, borrowed because the reader of this book already knows it. And a cell's kind --- an
ordinary term, a $\lambda$-term, an order-world command --- is decided by the engine and reported
in the reply; the notebook only shows the badge.

\section{Output references and forms}

Two conveniences borrowed from Mathematica. \texttt{\%} is the previous output, \texttt{\%\%} the
one before, \texttt{\%n} the output of evaluation $n$; the engine numbers every evaluation and
reports the number, the notebook shows it as \texttt{In[n]} and \texttt{Out[n]}, and the session
substitutes the referenced value before parsing goes any further, so the input interpretation shows
what \texttt{\%} stood for. And every output carries a small form chip: a matrix can be shown with
square or round brackets, as a bare grid, as a table, or in the engine's input syntax. The engine
prints once; the form is a typesetting choice the notebook keeps in the file.

\section{Exercises}

\begin{exercisebox}[A fourth host]
  The transport for a WebSocket is under thirty lines. Write it, and say what the engine has to
  change. (Nothing.)
\end{exercisebox}

\begin{exercisebox}[Rule 5 in practice]
  Add a field to \lc{EvaluateResult} that an old notebook would ignore and a new one would draw,
  and explain why the old notebook keeps working against the new engine and the new notebook
  against the old one.
\end{exercisebox}
